\documentclass[11pt]{article}
\usepackage[utf8]{inputenc}
\usepackage[T1]{fontenc}
\usepackage{amsmath,amssymb}
\usepackage{booktabs}
\usepackage{graphicx}
\usepackage{geometry}
\geometry{margin=1in}
\usepackage{authblk}
\usepackage[hidelinks]{hyperref}
\usepackage{caption}
\usepackage{xcolor}
\usepackage{titlesec}
\definecolor{navy}{HTML}{1A3A5C}
\definecolor{accent}{HTML}{2563EB}
\titleformat{\section}{\normalfont\Large\bfseries\color{navy}}{\thesection}{1em}{}
\titleformat{\subsection}{\normalfont\large\bfseries\color{accent}}{\thesubsection}{1em}{}

\title{\textbf{How Well Do Globally Trained Sentinel-1 Flood Mapping Methods\\
Transfer to Kenya? An Evaluation Against UNOSAT Rapid Mapping\\
of the 2024 Long-Rains Floods}}

\author[1]{Craig Carlos Ouma\thanks{%
Correspondence: \texttt{craigcarlos95@gmail.com}. Code and data:
\texttt{github.com/craigouma}}}
\affil[1]{Independent Researcher}
\date{\today}

\begin{document}
\maketitle

\begin{abstract}
\noindent
RESULTS-PENDING
\end{abstract}

\section{Introduction}
\label{sec:intro}

The March to May 2024 long rains were near or above average over most of
Kenya, and at several Nairobi stations, at Thika and at Machakos they were the
wettest season on record \cite{kmd2024}. Flooding followed along the Tana River
in Garissa and Tana River counties and across Nairobi and the neighbouring
counties of Kiambu and Machakos, and landslides struck Mai Mahiu and Kijabe
in Nakuru and Kiambu \cite{unosat2024}. By 18 June the National Disaster Operations Centre counted
315 people killed, 188 injured and 38 missing, and more than 293,200 people
displaced since 1 March \cite{ocha2024fu7}.

The United Nations Satellite Centre (UNOSAT) activated rapid mapping for the
event under code FL20240426KEN and published flood extents for Garissa and
the village of Mororo opposite it, for the Tana and Galana rivers, and for the
built-up area spanning Nairobi, Kiambu and Machakos \cite{unosat2024}. Those
maps were drawn from optical imagery, from Sentinel-2 and from very high
resolution Pléiades scenes, because optical imagery is what a human analyst
can interpret quickly. Optical sensors cannot see through cloud, and the long
rains are, by definition, the cloudiest part of the year. In the six days
either side of UNOSAT's reference dates, the three Sentinel-2 acquisitions over
Nairobi were 52\,\%, 89\,\% and 92\,\% cloud covered and the three over
Garissa 80\,\%, 36\,\% and 100\,\%, as reported for the whole 110\,km tile.

Sentinel-1 carries a C-band synthetic aperture radar (SAR) that images through
cloud, by day and night, and its data are free. Since the loss of Sentinel-1B in
December 2021, coverage has thinned: Garissa, for example, was imaged on a
single track every 12 days through 2024. Open water is dark in SAR because it
reflects the radar pulse away from the sensor, and this contrast is the basis
of every SAR flood mapping method.

The methods available off the shelf were developed or trained elsewhere.
Thresholding methods were tuned on European and South Asian floods, and the
most widely used labelled dataset for training deep learning models on
Sentinel-1, Sen1Floods11 \cite{bonafilia2020}, spans eleven flood events in
eleven countries, three of them African. Whether these methods hold up on
Kenyan landscapes, where semi-arid floodplains, cropland, informal
settlements and dense urban fabric each scatter radar differently, has not
been measured against an independent reference. This paper measures it.

The contributions are four. First, classical and deep learning Sentinel-1
methods are run over the 2024 event and scored against UNOSAT's reference
maps, separately for riverine and urban settings and broken down by land
cover, with uncertainty from a spatial block bootstrap. Second, the
sensitivity of that agreement to the number of days between the radar and
optical acquisitions is measured directly, by scoring every Sentinel-1 pass
within 24 days of each reference. Third, building exposure estimated from each
method is compared with UNOSAT's counts of affected structures. Fourth, the
method that performs best in riverine settings is applied to every pass over
the lower Tana between April and June 2024, as an application with its
measured error attached.

The paper is the companion to an earlier study in this series on satellite
indicators of drought and food insecurity in northern Kenya \cite{ouma2026drought},
and covers the opposite hydrological hazard with a second sensor type.

\section{Background}
\label{sec:background}

\subsection{SAR flood mapping}

The simplest method classifies a single post-event image by a backscatter
threshold. Because the threshold depends on sensor, incidence angle, wind and
land cover, it is usually fitted per scene, most often with Otsu's method
\cite{otsu1979}, which chooses the threshold that maximises the separation of
a bimodal histogram. A histogram over a whole scene is rarely bimodal when
water covers a small fraction of it, so split-based approaches fit the
threshold only on tiles that contain both classes
\cite{martinis2009,chini2017}; the same idea underlies fully automated
processing chains for Sentinel-1 \cite{twele2016,martinis2015}.

Change detection compares the flood image with a reference image from dry
conditions, usually as a log ratio, and labels as flooded the pixels that
became darker \cite{bazi2005,hostache2012}. It removes permanently dark
surfaces such as tarmac, sand and smooth bare soil, which single-image
thresholds confuse with water, at the cost of needing a suitable reference
image from the same orbit \cite{hostache2012}. Terrain descriptors such as
Height Above the Nearest Drainage (HAND) \cite{renno2008,nobre2011} are used
to remove water-like pixels on slopes and hilltops where flooding is
implausible.

The Copernicus Emergency Management Service now runs a global flood
monitoring product on every Sentinel-1 acquisition, as an ensemble of three
independently developed algorithms \cite{salamon2021,krullikowski2023}.

\subsection{Sen1Floods11}

Sen1Floods11 \cite{bonafilia2020} provides 4,831 chips of 512 by 512 pixels at
10\,m from eleven flood events, each with Sentinel-1 VV and VH backscatter
processed in Google Earth Engine. 446 chips carry hand-checked labels and are
split into training, validation and test sets, with the Bolivia event held out
as an unseen test; the remaining 4,385 carry weak labels from automated
classifications of Sentinel-1 or of coincident Sentinel-2. In the original
paper the best flood-water results came from a fully convolutional network
trained on the Sentinel-2-derived weak labels, with a mean per-chip IoU of
0.339 for flood water and 0.408 for all water on the test set, against 0.285
and 0.359 for Otsu thresholding of VH. Later work has reported Sentinel-1-only
results on the same benchmark with other architectures and loss functions
\cite{bai2021,konapala2021}; the highest published Sen1Floods11 numbers come
from optical foundation models \cite{jakubik2023,szwarcman2024}, which do not
use the radar channels.

\subsection{The urban SAR problem}

In built-up areas, walls and the ground form corner reflectors, and floodwater
between buildings strengthens this double-bounce return instead of darkening
it \cite{mason2012,mason2014}. Radar shadow and layover behind tall structures
add dark and bright artefacts unrelated to water. Floodwater in towns is
therefore often invisible to backscatter thresholds, and approaches that work
there rely on interferometric coherence or on combining SAR with an elevation
model of the urban surface \cite{mason2021}. The Sen1Floods11 validation data
include no urban flood events \cite{bonafilia2020}.

\section{Data}
\label{sec:data}

RESULTS-PENDING-DATA

\section{Methods}
\label{sec:methods}

RESULTS-PENDING-METHODS

\section{Results: Sen1Floods11 test performance}
\label{sec:s1f}

RESULTS-PENDING

\section{Results: transfer to Kenya}
\label{sec:kenya}

RESULTS-PENDING

\section{Exposure estimates against UNOSAT counts}
\label{sec:exposure}

RESULTS-PENDING

\section{Lower Tana flood time series}
\label{sec:tana}

RESULTS-PENDING

\section{Implications for flood monitoring in Kenya}
\label{sec:implications}

RESULTS-PENDING

\section{Limitations}
\label{sec:limitations}

RESULTS-PENDING

\section*{Reproducibility statement}

RESULTS-PENDING

\begin{thebibliography}{99}

\bibitem{bai2021}
Bai, Y., Wu, W., Yang, Z., Yu, J., Zhao, B., Liu, X., Yang, H., Mas, E., and
Koshimura, S. (2021). Enhancement of detecting permanent water and temporary
water in flood disasters by fusing Sentinel-1 and Sentinel-2 imagery using
deep learning algorithms: demonstration of Sen1Floods11 benchmark datasets.
\textit{Remote Sensing}, 13(11), 2220.

\bibitem{bazi2005}
Bazi, Y., Bruzzone, L., and Melgani, F. (2005). An unsupervised approach based
on the generalized Gaussian model to automatic change detection in
multitemporal SAR images. \textit{IEEE Transactions on Geoscience and Remote
Sensing}, 43(4), 874-887.

\bibitem{bonafilia2020}
Bonafilia, D., Tellman, B., Anderson, T., and Issenberg, E. (2020).
Sen1Floods11: a georeferenced dataset to train and test deep learning flood
algorithms for Sentinel-1. In \textit{2020 IEEE/CVF Conference on Computer
Vision and Pattern Recognition Workshops (CVPRW)}, pp. 835-845.

\bibitem{chini2017}
Chini, M., Hostache, R., Giustarini, L., and Matgen, P. (2017). A hierarchical
split-based approach for parametric thresholding of SAR images: flood
inundation as a test case. \textit{IEEE Transactions on Geoscience and Remote
Sensing}, 55(12), 6975-6988.

\bibitem{copdem}
European Space Agency and Airbus (2022). \textit{Copernicus DEM}.
\url{https://doi.org/10.5270/ESA-c5d3d65}

\bibitem{hostache2012}
Hostache, R., Matgen, P., and Wagner, W. (2012). Change detection approaches
for flood extent mapping: how to select the most adequate reference image from
online archives? \textit{International Journal of Applied Earth Observation
and Geoinformation}, 19, 205-213.

\bibitem{jakubik2023}
Jakubik, J., Roy, S., Phillips, C. E., Fraccaro, P., Godwin, D., Zadrozny, B.,
Szwarcman, D., Gomes, C., Nyirjesy, G., Edwards, B., Kimura, D., Simumba, N.,
Chu, L., Mukkavilli, S. K., Lambhate, D., Das, K., Bangalore, R., Oliveira,
D., Muszynski, M., Ankur, K., Ramasubramanian, M., Gurung, I., Khallaghi, S.,
Li, H., Cecil, M., Ahmadi, M., Kordi, F., Alemohammad, H., Maskey, M., Ganti,
R., Weldemariam, K., and Ramachandran, R. (2023). Foundation models for
generalist geospatial artificial intelligence. arXiv:2310.18660.

\bibitem{kmd2024}
Kenya Meteorological Department (2024). Climate outlook for the June-July-August
2024 season and rainfall review of the March-April-May 2024 ``long rains''
season. Ref. KMD/FCST/5-2024/SO/02, 31 May 2024.

\bibitem{konapala2021}
Konapala, G., Kumar, S. V., and Ahmad, S. K. (2021). Exploring Sentinel-1 and
Sentinel-2 diversity for flood inundation mapping using deep learning.
\textit{ISPRS Journal of Photogrammetry and Remote Sensing}, 180, 163-173.

\bibitem{krullikowski2023}
Krullikowski, C., Chow, C., Wieland, M., Martinis, S., Bauer-Marschallinger,
B., Roth, F., Matgen, P., Chini, M., Hostache, R., Li, Y., and Salamon, P.
(2023). Estimating ensemble likelihoods for the Sentinel-1-based global flood
monitoring product of the Copernicus Emergency Management Service.
\textit{IEEE Journal of Selected Topics in Applied Earth Observations and
Remote Sensing}, 16, 6917-6930.

\bibitem{martinis2009}
Martinis, S., Twele, A., and Voigt, S. (2009). Towards operational near
real-time flood detection using a split-based automatic thresholding procedure
on high resolution TerraSAR-X data. \textit{Natural Hazards and Earth System
Sciences}, 9(2), 303-314.

\bibitem{martinis2015}
Martinis, S., Kersten, J., and Twele, A. (2015). A fully automated TerraSAR-X
based flood service. \textit{ISPRS Journal of Photogrammetry and Remote
Sensing}, 104, 203-212.

\bibitem{mason2012}
Mason, D. C., Davenport, I. J., Neal, J. C., Schumann, G. J.-P., and Bates,
P. D. (2012). Near real-time flood detection in urban and rural areas using
high-resolution synthetic aperture radar images. \textit{IEEE Transactions on
Geoscience and Remote Sensing}, 50(8), 3041-3052.

\bibitem{mason2014}
Mason, D. C., Giustarini, L., Garcia-Pintado, J., and Cloke, H. L. (2014).
Detection of flooded urban areas in high resolution synthetic aperture radar
images using double scattering. \textit{International Journal of Applied Earth
Observation and Geoinformation}, 28, 150-159.

\bibitem{mason2021}
Mason, D. C., Dance, S. L., and Cloke, H. L. (2021). Floodwater detection in
urban areas using Sentinel-1 and WorldDEM data. \textit{Journal of Applied
Remote Sensing}, 15(3), 032003.

\bibitem{nemni2020}
Nemni, E., Bullock, J., Belabbes, S., and Bromley, L. (2020). Fully
convolutional neural network for rapid flood segmentation in synthetic
aperture radar imagery. \textit{Remote Sensing}, 12(16), 2532.

\bibitem{nobre2011}
Nobre, A. D., Cuartas, L. A., Hodnett, M., Renn\'o, C. D., Rodrigues, G.,
Silveira, A., Waterloo, M., and Saleska, S. (2011). Height Above the Nearest
Drainage: a hydrologically relevant new terrain model. \textit{Journal of
Hydrology}, 404(1-2), 13-29.

\bibitem{ocha2024fu7}
UN Office for the Coordination of Humanitarian Affairs (2024). Kenya: Heavy
Rains and Flooding Update, Flash Update \#7 (19 June 2024).
\url{https://www.unocha.org/publications/report/kenya/kenya-heavy-rains-and-flooding-update-flash-update-7-19-june-2024}

\bibitem{otsu1979}
Otsu, N. (1979). A threshold selection method from gray-level histograms.
\textit{IEEE Transactions on Systems, Man, and Cybernetics}, 9(1), 62-66.

\bibitem{ouma2026drought}
Ouma, C. C. (2026). Satellite-derived drought indicators as early predictors of
food insecurity risk in northern Kenya's arid and semi-arid lands. Preprint.
\url{https://craigouma.github.io/eo-drought-food-security-kenya/}

\bibitem{pekel2016}
Pekel, J.-F., Cottam, A., Gorelick, N., and Belward, A. S. (2016).
High-resolution mapping of global surface water and its long-term changes.
\textit{Nature}, 540(7633), 418-422.

\bibitem{planetary2022}
Microsoft Open Source, McFarland, M., Emanuele, R., Morris, D., and
Augspurger, T. (2022). \textit{microsoft/PlanetaryComputer: October 2022}.
Zenodo. \url{https://doi.org/10.5281/zenodo.7261897}

\bibitem{renno2008}
Renn\'o, C. D., Nobre, A. D., Cuartas, L. A., Soares, J. V., Hodnett, M. G.,
Tomasella, J., and Waterloo, M. J. (2008). HAND, a new terrain descriptor using
SRTM-DEM: mapping terra-firme rainforest environments in Amazonia.
\textit{Remote Sensing of Environment}, 112(9), 3469-3481.

\bibitem{ronneberger2015}
Ronneberger, O., Fischer, P., and Brox, T. (2015). U-Net: convolutional
networks for biomedical image segmentation. In \textit{Medical Image Computing
and Computer-Assisted Intervention, MICCAI 2015}, pp. 234-241. Springer.

\bibitem{salamon2021}
Salamon, P., McCormick, N., Reimer, C., Clarke, T., Bauer-Marschallinger, B.,
Wagner, W., Martinis, S., Chow, C., B\"ohnke, C., Matgen, P., Chini, M.,
Hostache, R., Molini, L., Fiori, E., and Walli, A. (2021). The new, systematic
global flood monitoring product of the Copernicus Emergency Management
Service. In \textit{2021 IEEE International Geoscience and Remote Sensing
Symposium (IGARSS)}, pp. 1053-1056.

\bibitem{sirko2021}
Sirko, W., Kashubin, S., Ritter, M., Annkah, A., Bouchareb, Y. S. E., Dauphin,
Y., Keysers, D., Neumann, M., Cisse, M., and Quinn, J. (2021).
Continental-scale building detection from high resolution satellite imagery.
arXiv:2107.12283.

\bibitem{szwarcman2024}
Szwarcman, D., Roy, S., Fraccaro, P., et al. (2024). Prithvi-EO-2.0: a
versatile multi-temporal foundation model for Earth observation applications.
arXiv:2412.02732.

\bibitem{twele2016}
Twele, A., Cao, W., Plank, S., and Martinis, S. (2016). Sentinel-1-based flood
mapping: a fully automated processing chain. \textit{International Journal of
Remote Sensing}, 37(13), 2990-3004.

\bibitem{unosat2024}
UNITAR Satellite Centre (UNOSAT) (2024). Flood products for activation
FL20240426KEN (GDACS 1102561), Kenya, 29 April to 5 May 2024. Humanitarian
Data Exchange. \url{https://data.humdata.org/organization/unosat}

\bibitem{zanaga2022}
Zanaga, D., Van De Kerchove, R., Daems, D., De Keersmaecker, W., Brockmann, C.,
Kirches, G., Wevers, J., Cartus, O., Santoro, M., Fritz, S., Lesiv, M.,
Herold, M., Tsendbazar, N.-E., Xu, P., Ramoino, F., and Arino, O. (2022).
\textit{ESA WorldCover 10 m 2021 v200}. Zenodo.
\url{https://doi.org/10.5281/zenodo.7254221}

\end{thebibliography}

\end{document}
